\documentclass[journal]{IEEEtran}

\usepackage{amsmath}
\usepackage{amssymb}
\usepackage{graphicx}
\usepackage{booktabs}
\usepackage{algorithm}
\usepackage{algorithmic}
\usepackage{cite}
\usepackage{url}
\usepackage{hyperref}

\hypersetup{
  colorlinks=true,
  linkcolor=black,
  citecolor=black,
  urlcolor=blue
}

\graphicspath{{figures/}}

\begin{document}

\title{When Does Gradient Refinement Help Sampling-Based MPC?\\
Localizing the Benefit of Diff-MPPI to Contact-Rich Manipulation}

\author{Ryohei~Sasaki%
\thanks{Manuscript received XXX; revised XXX.}%
}

\markboth{IEEE Robotics and Automation Letters}%
{Sasaki: Diff-MPPI}

\maketitle

% =============================================================================
\begin{abstract}
Sampling-based model predictive controllers such as MPPI handle nonlinear
dynamics and nonconvex costs effectively, and recent work shows MPPI is
equivalent to a preconditioned gradient descent step. A natural question is
whether \emph{adding} explicit gradient steps to MPPI helps, and if so, where.
We study a deliberately minimal hybrid controller---Diff-MPPI---that appends a
short forward-mode autodiff refinement (a few local gradient steps) to the
standard MPPI update, parallelized on GPU so it preserves most of the sampling
budget at matched wall-clock time. Our central finding is a \emph{localization}
rather than a blanket win. On smooth dynamic-obstacle navigation and 7-DOF
serial-arm reaching, a literature-faithful One-Step CDF-MPPI baseline,
reimplemented in the same harness, matches or beats the entire Diff-MPPI family
at $8$--$47\times$ lower per-step compute; and under a fair shared-goal-configuration
protocol the gradient refinement does \emph{not} outperform vanilla MPPI. The
benefit appears specifically in \emph{contact-rich} dynamics. On a
differentiable planar pushing task where the model gradient flows through a
smooth contact model, Diff-MPPI pushes a box to a target pose with success
$1.00$ while vanilla MPPI reaches $0.00$ \emph{even at $16\times$ the samples}
(a cheaper per-step budget), with success monotone in the number of gradient
steps; a second contact task and a disk-pushing matched-wall-clock efficiency
result corroborate the effect, and the advantage persists under $\pm$30--40\%
contact-model mismatch. The contribution is thus a clearly scoped,
mechanism-supported positive result with a rigorous negative control: gradient
refinement buys what sampling cannot precisely when information is carried by the
contact gradient.
\end{abstract}

\begin{IEEEkeywords}
Motion planning, model predictive control, sampling-based optimization,
gradient-based refinement, differentiable simulation, non-prehensile
manipulation, GPU-accelerated control
\end{IEEEkeywords}

% =============================================================================
\section{Introduction}
\label{sec:introduction}

Model Predictive Path Integral (MPPI) control is widely used for real-time
nonlinear trajectory optimization because it handles arbitrary cost functions
through parallel rollout sampling on GPU~\cite{williams2017mppi}. Recent
theoretical work by Fazlyab et al.~\cite{fazlyab2026mppi_gd} shows that the
MPPI update is exactly a preconditioned gradient descent step with unit step
size on a KL-regularized objective. This invites an obvious construction: after
the MPPI weighted update, take a few \emph{explicit} gradient steps on the
nominal control sequence using automatic differentiation. We call this minimal
hybrid Diff-MPPI. The question this paper answers is not ``can we build it'' but
``\emph{when does it actually help?}''---and, just as importantly, when it does
not.

We answer with a deliberately adversarial evaluation. Rather than reporting only
the tasks where the hybrid wins, we (i) implement a strong, literature-faithful
direct baseline---One-Step CDF-MPPI~\cite{cdf_mppi}---inside the same CUDA
harness, and (ii) equalize the most important confound (goal specification) so
that any remaining difference is attributable to the controller. Under this
protocol, the picture on \emph{smooth} tasks is humbling: on 2D dynamic-obstacle
navigation and 7-DOF reaching, CDF-MPPI matches or beats the Diff-MPPI family at
a fraction of the per-step compute, and once every method is handed the same
goal configuration the gradient refinement is statistically indistinguishable
from vanilla MPPI. A short autodiff correction does not, by itself, beat strong
sampling on smooth dynamics.

The benefit reappears---sharply---when the dynamics gradient carries information
that sampling cannot cheaply discover: \emph{contact}. We build a differentiable
planar pushing task with a smooth (softplus / signed-distance) contact model and
forward-mode dual-number autodiff that flows through contact. There, pushing a
box to a target pose, Diff-MPPI succeeds $1.00$ of the time while vanilla MPPI
succeeds $0.00$ of the time even at $16\times$ the sample budget, and success
rises monotonically with the number of gradient steps. The effect reproduces on
a second, mechanically distinct contact task and as a matched-wall-clock
efficiency gain on disk pushing.

Our contributions are:
\begin{enumerate}
  \item \textbf{A minimal, training-free hybrid MPPI} that refines the
    post-sampling control mean with forward-mode autodiff, including through a
    smooth contact model, parallelized to preserve most of the sampling budget
    at matched wall-clock time (Sec.~\ref{sec:method}).
  \item \textbf{A rigorous negative control.} A literature-faithful CDF-MPPI
    baseline in the same harness, plus a fair shared-goal-configuration
    protocol, showing that on smooth dynamic-navigation and 7-DOF reaching the
    gradient refinement provides \emph{no} advantage over vanilla MPPI and is
    dominated on compute by CDF-MPPI (Sec.~\ref{sec:results_smooth}).
  \item \textbf{A localized positive result.} On differentiable-contact pushing,
    a categorical success-rate win (box pose: $1.00$ vs $0.00$ at $16\times$
    samples) that is monotone in gradient steps, corroborated by a second
    contact task and a disk-pushing matched-time efficiency gain
    (Sec.~\ref{sec:results_contact}). This localizes the benefit of gradient
    refinement in sampling-based MPC to contact-rich dynamics.
\end{enumerate}

% =============================================================================
\section{Related Work}
\label{sec:related}

\textbf{MPPI and path-integral control.}
MPPI~\cite{williams2017mppi} performs weighted averaging of sampled trajectories
using the exponential cost-to-go as importance weights. Fazlyab et
al.~\cite{fazlyab2026mppi_gd} show this update is a preconditioned gradient
descent step, connecting MPPI to optimization-based MPC and directly motivating
the question of whether explicit gradient steps add anything.

\textbf{Sampling + gradient hybrids.}
CEM-GD~\cite{bharadhwaj2020cemgd} combines cross-entropy sampling with gradient
descent for model-based RL planning. MPPI-IPDDP~\cite{mppi_ipddp} smooths MPPI
trajectories via DDP inside a convex corridor. Our refinement is simpler---pure
autodiff on the post-MPPI control sequence---but our emphasis is different: we
characterize \emph{where} such a hybrid helps and provide a negative control on
smooth tasks rather than asserting a general improvement.

\textbf{Direct sampling-based baseline.}
One-Step CDF-MPPI~\cite{cdf_mppi} shapes a single-step MPPI cost using the angle
between sampled motion and the ascent direction of a configuration-space
distance field, and is evaluated on 7-DOF reaching---the same regime as our
smooth manipulation tasks. We reimplement it faithfully (analytic
margin-derived distance field as primary, a neural variant as an ablation) as a
direct, mechanistically distinct baseline.

\textbf{Feedback-enhanced and learned-sampling MPPI.}
Feedback-MPPI~\cite{belvedere2026feedback_mppi} computes local feedback gains
from rollout sensitivity; Biased-MPPI~\cite{trevisan2024biased_mppi} uses
ancillary controllers; Step-MPPI~\cite{step_mppi} learns a sampling
distribution. These modify sampling or feedback; our method instead adds an
explicit gradient step, and we find it complementary only in contact-rich
settings.

\textbf{Differentiable simulation and contact.}
Differentiable physics provides analytic gradients through contact for
manipulation and locomotion~\cite{diffmpc2025,cunrto2026}. Non-prehensile
pushing is a canonical contact-rich task whose dynamics are governed by frictional
contact wrenches. We use a smooth (softplus penetration, signed-distance) contact
model so that forward-mode autodiff yields an exact gradient of the rollout cost
through contact, which is precisely the signal our positive result exploits.

% =============================================================================
\section{Method}
\label{sec:method}

\subsection{Controller Structure}
\label{sec:controller_structure}

At each control step, Diff-MPPI executes three phases
(Algorithm~\ref{alg:diff_mppi}): a standard MPPI weighted update, a parallel
cost-gradient evaluation, and $N_g$ local gradient-descent steps on the nominal
control sequence with per-timestep norm clipping. No component is learned and no
auxiliary structure beyond the cost gradient is built.

\begin{algorithm}[t]
\caption{Diff-MPPI Controller (one control step)}
\label{alg:diff_mppi}
\begin{algorithmic}[1]
\REQUIRE Current state $\mathbf{x}_0$, nominal controls $\mathbf{u}_{0:T-1}$, sample count $K$, gradient steps $N_g$, step size $\alpha$, clip threshold $c$
\ENSURE Updated nominal controls $\mathbf{u}_{0:T-1}$, applied control $\mathbf{u}_0$
\STATE \textbf{--- Phase 1: MPPI Sampling Update ---}
\STATE Sample $K$ perturbations $\boldsymbol{\epsilon}^{(k)} \sim \mathcal{N}(0, \Sigma)$
\STATE Roll out $K$ trajectories: $\tilde{\mathbf{u}}^{(k)} = \mathbf{u} + \boldsymbol{\epsilon}^{(k)}$
\STATE Compute trajectory costs $S^{(k)}$ \COMMENT{$K$ GPU threads}
\STATE Compute weights $w^{(k)} \propto \exp(-\frac{1}{\lambda} S^{(k)})$
\STATE $\mathbf{u} \leftarrow \mathbf{u} + \sum_k w^{(k)} \boldsymbol{\epsilon}^{(k)}$
\STATE \textbf{--- Phase 2: Parallel Cost Gradient (autodiff) ---}
\STATE For each control parameter, evaluate $\partial S / \partial \mathbf{u}_t$ by a forward-mode dual-number rollout \COMMENT{parallel over parameters}
\STATE \textbf{--- Phase 3: Local Gradient Refinement ---}
\FOR{$i = 1$ \TO $N_g$}
  \STATE Clip: $\mathbf{g}_t \leftarrow \mathbf{g}_t \cdot \min(1,\; c / \|\mathbf{g}_t\|)$ for all $t$
  \STATE $\mathbf{u}_t \leftarrow \mathbf{u}_t - \alpha\, \mathbf{g}_t$ for all $t$
  \STATE Recompute $\mathbf{g}$ by a fresh dual-number rollout
\ENDFOR
\RETURN $\mathbf{u}_{0:T-1}$, apply $\mathbf{u}_0$
\end{algorithmic}
\end{algorithm}

\subsection{Forward-Mode Autodiff Through Smooth Contact}
\label{sec:autodiff_contact}

The gradient $\partial S/\partial \mathbf{u}$ is obtained by a forward-mode
dual-number rollout: every state and intermediate quantity carries a value and a
derivative with respect to the active control parameter, so a single rollout
returns both the cost and its exact gradient. The key requirement for the
contact tasks is that the contact model be \emph{smooth}. For planar pushing we
use a softplus penetration $\mathrm{pen} = \mathrm{softplus}(r - \mathrm{sd})$ of
the pusher into the object signed distance $\mathrm{sd}$, a normal force
proportional to penetration, and---for the box---a torque from the off-centre
contact wrench computed from the box signed-distance gradient. Because every
operation (softplus, signed distance of a box, the contact normal, the wrench)
is differentiable, the dual-number rollout differentiates the entire
contact-mediated trajectory, giving the refinement a usable gradient exactly
where pure sampling struggles.

\subsection{Compute Budget Analysis}
\label{sec:budget}

After parallelization, the refinement adds a bounded overhead per step. The
comparisons in Sec.~\ref{sec:results} therefore use \emph{matched wall-clock per
control step} (an identical timing meter wrapped around the controller update in
all binaries) as the primary budget axis; sample count and replan rate are
disclosed as secondary columns. Crucially, in the contact result the hybrid is
compared against vanilla MPPI given \emph{more} samples at \emph{lower} per-step
cost, so the gradient cannot be ``buying'' quality by escaping the budget.

% =============================================================================
\section{Experimental Setup}
\label{sec:experiments}

\subsection{Smooth Tasks (negative-control domains)}

\textbf{2D dynamic-obstacle navigation.} A bicycle model with
\texttt{dynamic\_crossing} (one crossing obstacle) and \texttt{dynamic\_slalom}
(multiple moving obstacles).

\textbf{7-DOF serial-arm reaching.} A Panda-like 7-DOF arm with 14D state, 7D
control, 3D workspace obstacles, and analytical kinematics. Scenarios:
\texttt{7dof\_shelf\_reach} (static) and \texttt{7dof\_dynamic\_avoid} (moving
3D obstacle).

\subsection{Contact-Rich Tasks (positive-result domains)}

\textbf{Differentiable planar pushing.} A quasi-static point pusher acts on an
object through the smooth contact model of Sec.~\ref{sec:autodiff_contact}.
\texttt{push\_straight}/\texttt{push\_diagonal} push a disk to a target
position; \texttt{box\_align} and \texttt{box\_pivot} push a rectangular box to
a target \emph{pose} $(x,y,\theta)$, which requires off-centre contact to
generate torque. \texttt{box\_pivot} is mechanically distinct from
\texttt{box\_align} (opposite rotation handedness, a different contact face, and
a tighter angular tolerance). \texttt{box\_turn} (long translate~+~rotate) is
included as a stress case.

\subsection{Baselines}

\textbf{CDF-MPPI}~\cite{cdf_mppi}: a literature-faithful One-Step CDF-MPPI in the
same harness, using an analytic margin-derived configuration-space distance
field (a neural distance field is reported as an ablation). \textbf{Vanilla
MPPI} and a family of in-harness \texttt{feedback\_mppi} variants
(sensitivity-, covariance-, and release-style gains) are also included.
Diff-MPPI-$N_g$ denotes the hybrid with $N_g$ gradient steps.

\subsection{Fair Goal-Configuration Protocol}

CDF-MPPI is natively handed a collision-free goal \emph{configuration}
$\mathbf{q}_{\text{goal}}$ (effectively a solved IK), whereas the Diff-MPPI
family default sees only a workspace goal. To remove this confound we add a
gated joint-space goal cost to the Diff-MPPI binary and feed it the same
$\mathbf{q}_{\text{goal}}$ (default weight $0$, so published numbers reproduce
byte-identically). All smooth-task comparisons are reported under this shared
condition.

% =============================================================================
\section{Results}
\label{sec:results}

\subsection{Where it does \emph{not} help: smooth tasks}
\label{sec:results_smooth}

On \texttt{7dof\_shelf\_reach} under matched wall-clock
(Table~\ref{tab:cdf_reach}), the literature-faithful CDF-MPPI reaches the goal
collision-free at $0.05$\,ms/step where the best Diff-MPPI variant manages only
$0.50$ success at roughly $9\times$ the per-step compute. The naive gap is
partly CDF-MPPI's free goal configuration; the fair rematch
(Table~\ref{tab:fair}) closes most of it---and, decisively, once every method is
handed the same $\mathbf{q}_{\text{goal}}$, \texttt{diff\_mppi\_3} ($0.62$) does
\emph{not} beat vanilla \texttt{mppi} ($0.88$). The same pattern holds on the
dynamic 7-DOF task, where a reactive CDF-MPPI is the most reliable and cheapest
controller ($1.00$ success, zero collisions, $0.24$\,ms/step) while the
Diff-MPPI family tops out near $0.88$, matching vanilla MPPI. Across these smooth
tasks the autodiff refinement shows no demonstrable advantage; this is the
negative control that scopes the positive result below.

\begin{table}[t]
\centering
\caption{%
  \texttt{7dof\_shelf\_reach}, matched wall-clock, 4 seeds. CDF-MPPI is the
  literature-faithful baseline; the Diff-MPPI family does not beat it on success
  or compute.%
}
\label{tab:cdf_reach}
\begin{tabular}{lccc}
\toprule
Planner & Success & Final Dist & ms/step \\
\midrule
cdf\_mppi (analytic)   & \textbf{1.00} & 0.144 & \textbf{0.05} \\
diff\_mppi\_1          & 0.50 & 0.274 & 0.46 \\
diff\_mppi\_3          & 0.25 & 0.328 & 0.75 \\
mppi                   & 0.25 & 0.340 & 0.38 \\
feedback\_mppi\_ref    & 0.00 & 0.398 & 1.66 \\
\bottomrule
\end{tabular}
\end{table}

\begin{table}[t]
\centering
\caption{%
  Fair rematch on \texttt{7dof\_shelf\_reach} (8 seeds): with the \emph{same}
  goal configuration, every method improves and \texttt{diff\_mppi\_3} does not
  exceed vanilla \texttt{mppi}.%
}
\label{tab:fair}
\begin{tabular}{lccc}
\toprule
Planner & Success (workspace) & Success (shared config) & ms/step \\
\midrule
mppi                & 0.38 & \textbf{0.88} & 0.39 \\
diff\_mppi\_1       & 0.50 & 0.62 & 0.54 \\
diff\_mppi\_3       & 0.25 & 0.62 & 0.85 \\
feedback\_mppi\_ref & 0.00 & 1.00 & 2.33 \\
\bottomrule
\end{tabular}
\end{table}

\begin{figure}[t]
  \centering
  \includegraphics[width=\columnwidth]{fig_cdf_vs_diff}
  \caption{%
    Compute--quality on the smooth \texttt{7dof\_shelf\_reach} task. The
    literature-faithful CDF-MPPI sits at the top-left (full success at
    $0.05$\,ms/step); the Diff-MPPI family does not dominate it, and vanilla
    \texttt{mppi} matches \texttt{diff\_mppi\_3} under the fair shared-goal
    condition. This is the negative control that scopes the contact result.%
  }
  \label{fig:cdf_vs_diff}
\end{figure}

\subsection{Where it helps: contact-rich pushing}
\label{sec:results_contact}

\textbf{Box pose alignment (categorical win).} On \texttt{box\_align}, pushing a
box to a target pose requires torque from off-centre contact, and pure sampling
plateaus on the fine orientation correction. Table~\ref{tab:box_align} shows the
headline: \texttt{diff\_mppi\_5} reaches success $1.00$ (at $K\!=\!1024$) while
\textbf{vanilla MPPI reaches $0.00$ even at $K\!=\!4096$}---$16\times$ the
samples, at a \emph{lower} per-step cost than the hybrid---and success is
monotone in the number of gradient steps ($0.00 \!\to\! 0.50 \!\to\! 1.00$ for
$N_g\!=\!1,3,5$). No sample budget buys what the gradient provides.

\begin{table}[t]
\centering
\caption{%
  \texttt{box\_align} differentiable box-pose pushing, 8 seeds. Vanilla MPPI
  fails at every budget, including $16\times$ samples; success is monotone in
  gradient steps. \textbf{Bold}: proposed method.%
}
\label{tab:box_align}
\begin{tabular}{lccc}
\toprule
Planner & Success ($K\!=\!256$) & Success ($K\!=\!1024$) & ms/step \\
\midrule
mppi ($K\!\le\!4096$)        & 0.00 & 0.00 & 0.13--0.52 \\
diff\_mppi\_1                & 0.12 & 0.00 & 0.85 \\
diff\_mppi\_3                & 0.25 & 0.50 & 2.20 \\
\textbf{diff\_mppi\_5}       & \textbf{0.62} & \textbf{1.00} & 3.08 \\
\bottomrule
\end{tabular}
\end{table}

\begin{figure}[t]
  \centering
  \includegraphics[width=\columnwidth]{fig_contact_filmstrip}
  \caption{%
    Executed box poses on \texttt{box\_align} (target pose: green dashed; pusher
    path: grey). Vanilla MPPI (left) plateaus with the box misaligned and never
    reaches the pose ($240$ steps); \texttt{diff\_mppi\_5} (right) rotates the box
    into the target pose in $47$ steps. The gradient through contact supplies the
    fine orientation correction that sampling cannot.%
  }
  \label{fig:filmstrip}
\end{figure}

\begin{figure}[t]
  \centering
  \includegraphics[width=\columnwidth]{fig_box_samples}
  \caption{%
    \texttt{box\_align} success vs.\ sample budget $K$. Vanilla MPPI stays at
    $0.00$ even at $K\!=\!4096$ ($16\times$ samples, a cheaper per-step budget
    than the hybrid), while \texttt{diff\_mppi\_5} succeeds. Samples cannot
    replace the gradient.%
  }
  \label{fig:box_samples}
\end{figure}

\textbf{Corroboration on distinct mechanics (\texttt{box\_pivot}).} To guard
against a single lucky scenario, \texttt{box\_pivot} uses the opposite rotation
handedness, a different contact face, and a tighter angular tolerance. The same
two signatures reproduce (Table~\ref{tab:box_pivot}): only the deepest-gradient
controller ever crosses the tolerance while vanilla MPPI succeeds $0.00$ at every
$K$, and the \emph{continuous} angular residual is strictly monotone in gradient
steps ($0.193 \!\to\! 0.139 \!\to\! 0.124 \!\to\! 0.112$) and $K$-independent. The
success rate here is deliberately modest---\texttt{box\_pivot} sits near the
rotation ceiling of a single fixed-face push---so we report it primarily on the
continuous monotone metric, with the binary win as corroboration rather than a
threshold-tuned headline.

\begin{table}[t]
\centering
\caption{%
  \texttt{box\_pivot} corroboration, 8 seeds, $K\!=\!1024$. The continuous
  angular residual (lower is better) is strictly monotone in gradient steps;
  only \texttt{diff\_mppi\_5} crosses the tolerance.%
}
\label{tab:box_pivot}
\begin{tabular}{lcc}
\toprule
Planner & Success & Angular residual \\
\midrule
mppi          & 0.00 & 0.193 \\
diff\_mppi\_1 & 0.00 & 0.139 \\
diff\_mppi\_3 & 0.00 & 0.124 \\
\textbf{diff\_mppi\_5} & \textbf{0.50} & \textbf{0.112} \\
\bottomrule
\end{tabular}
\end{table}

\textbf{Robustness to contact-model mismatch.} A natural objection is that the
gradient win is an artifact of a contact model whose parameters exactly match the
plant. We test this directly: the \emph{true} plant's contact mobility (force and
torque gains) is scaled by a factor $G$ while the controller's rollout and
gradient keep the nominal gains---its contact model is deliberately wrong by $G$
($G\!=\!1$ recovers the matched case). Table~\ref{tab:mismatch} shows the win is
not a matched-model knife's edge: across $G \in [0.7, 1.4]$ ($\pm 30$--$40\%$
contact-mobility error), \texttt{diff\_mppi\_5} retains a categorical advantage
over vanilla MPPI, which stays at $0.00$ throughout. The degradation is
asymmetric and interpretable: under-modelled mobility ($G\!<\!1$) is \emph{better}
than matched (the controller commits to sustained contact and converges
conservatively), while over-modelled mobility ($G\!>\!1$) overshoots the fine
orientation and degrades monotonically. The advantage lives in the stiff-contact
regime where sampling overshoots; at $G\!=\!0.6$ the dynamics become gentle enough
that vanilla MPPI also reaches $1.00$ and the task no longer discriminates.

\begin{table}[t]
\centering
\caption{%
  Contact-model mismatch on \texttt{box\_align}, 8 seeds. The plant's contact
  mobility is scaled by $G$ while the controller keeps nominal gains. The
  gradient win survives $\pm 30$--$40\%$ model error; vanilla MPPI fails at every
  $G$ in that band even at $16\times$ samples.%
}
\label{tab:mismatch}
\begin{tabular}{lcc}
\toprule
Plant/model gain $G$ & mppi ($K\!=\!4096$) & \textbf{diff\_mppi\_5} ($K\!=\!1024$) \\
\midrule
0.6 (task easy for all) & 1.00 & 1.00 \\
0.7  & 0.00 & \textbf{1.00} \\
0.85 & 0.00 & \textbf{1.00} \\
1.0 (matched)  & 0.00 & \textbf{0.75} \\
1.2  & 0.00 & \textbf{0.50} \\
1.4  & 0.00 & \textbf{0.25} \\
1.6  & 0.00 & 0.12 \\
\bottomrule
\end{tabular}
\end{table}

\textbf{Matched-wall-clock efficiency on disk pushing.} On \texttt{push\_straight},
where all methods eventually reach the goal, the discriminator is efficiency.
Vanilla MPPI saturates at $\sim$30 steps regardless of $K$ (tested $512$--$2048$),
while \texttt{diff\_mppi\_3} reaches the goal in $\sim$26 steps at matched
wall-clock---an edge that additional samples cannot erase. Together, three
contact tasks (disk efficiency, two box-pose wins) localize the Diff-MPPI benefit
to contact-rich dynamics.

\subsection{Mechanism and Boundary}
\label{sec:mechanism}

\begin{figure}[t]
  \centering
  \includegraphics[width=\columnwidth]{fig_contact_monotone}
  \caption{%
    The gradient is the active ingredient. (a) \texttt{box\_align} success and
    (b) \texttt{box\_pivot} continuous angular residual (axis inverted so better
    is up) both improve monotonically with the number of gradient steps $N_g$,
    with the sampler held fixed; $N_g\!=\!0$ is vanilla MPPI.%
  }
  \label{fig:monotone}
\end{figure}

The monotonicity of success (and of the continuous residual) in the number of
gradient steps (Fig.~\ref{fig:monotone}) is direct evidence that the autodiff
refinement, not the sampling, is the active ingredient in the contact wins:
holding the sampler fixed and adding gradient steps strictly improves the
executed pose. The boundary is
\texttt{box\_turn} (long translate~+~rotate), which \emph{no} method solves: its
bottleneck is translation (sample-friendly), not orientation, and a longer
planning horizon reduces the residual for the hybrid but vanilla MPPI keeps pace.
Single-point non-prehensile pushing fundamentally needs contact-point switching
for such a maneuver, which receding-horizon mean refinement does not plan. We
report \texttt{box\_turn} as a stated limit rather than tuning a win on the wrong
axis.

% Figures for the contact-rich framing are generated by
% scripts/plot_contact_figures.py (data dumped via
% `bin/benchmark_diff_mppi_pushing_box --dump-traj build/box_traj`):
%   fig_contact_filmstrip, fig_box_samples, fig_contact_monotone, fig_cdf_vs_diff.
% The previous dynamic-navigation figures (fig_pareto, fig_7dof, fig_mechanism,
% fig_scenarios) describe the superseded headline and are intentionally omitted.

% =============================================================================
\section{Limitations}
\label{sec:limitations}

The positive result uses a simplified planar quasi-static contact model rather
than a validated rigid-body simulator with friction; extending the gradient
through a higher-fidelity contact model is the natural next step. We partially
de-risk this with the contact-model mismatch study (Table~\ref{tab:mismatch}),
which shows the advantage survives $\pm 30$--$40\%$ error in the plant's contact
mobility, but a full sim-to-real validation remains future work. The contact
tasks are non-prehensile pushing only, and \texttt{box\_turn} marks a maneuver
the method does not yet solve. The CDF-MPPI baseline, while literature-faithful,
is an in-harness reimplementation, and the smooth-task models (bicycle, simplified
7-DOF) are custom rather than standardized suites. Finally, the contact study is
in simulation; sim-to-real transfer of the differentiable-contact gradient is
untested.

% =============================================================================
\section{Conclusion}
\label{sec:conclusion}

We asked when adding explicit gradient steps to MPPI helps. The honest answer is
narrow and, we argue, more useful than a blanket claim: on smooth
dynamic-navigation and 7-DOF reaching, a literature-faithful CDF-MPPI baseline
matches or beats the Diff-MPPI family at a fraction of the compute, and under a
fair goal-configuration protocol the refinement does not beat vanilla MPPI. But
when the dynamics gradient flows through contact, a few autodiff steps achieve
what no sample budget can: a categorical box-pose pushing win ($1.00$ vs $0.00$
at $16\times$ samples), monotone in gradient steps, reproduced on a second
contact task and as a matched-time efficiency gain. Diff-MPPI's contribution is
real but localized---contact-rich dynamics are where the model gradient carries
information that is expensive to sample.

% =============================================================================
\bibliographystyle{IEEEtran}
\bibliography{references}

\end{document}
